\input{../tex-inputs/latex-leseansicht-vorspann}

\section[Arthur Schnitzler an Gustav Schwarzkopf, 22. 4. 1927]{L04468 Arthur Schnitzler an Gustav Schwarzkopf, 22. 4. 1927}
\nopagebreak\mylabel{L04468v}
\rehead{ }\normalsize\beginnumbering\briefempfaengerindex{Schwarzkopf, Gustav@\textsc{Schwarzkopf, Gustav}!zzzSchnitzler, Arthur@\emph{von Arthur Schnitzler}!1927-04-221@{22. 4. 1927}|(be}
\toendnotes[C]{\smallbreak\pagebreak[2]}
\correspDesc{Versand  durch Arthur Schnitzler am 22. 4. 1927 in Venedig
\newline{}Übermittlung  am 22. 4. 1927 in Venedig
\newline{}Erhalt  durch Gustav Schwarzkopf im Zeitraum [23. 4. 1927 – 27. 4. 1927?] in Wien}\toendnotes[C]{\smallbreak}
\Standort{DLA, A:Schnitzler, HS.1985.1.1897.}
\physDesc{Bildpostkarte, 106 Zeichen
\newline{}Handschrift: Bleistift, lateinische Kurrent
\newline{}Versand: Stempel: »\nobreak{}\oindex{Venedig@\textbf{Venedig}|pwk}Venezia, 22. IV. 1927, 22–23, Ferrovia\nobreak{}«.  }\pstart{}{\pb}Austria\oindex{Österreich@\textbf{Österreich}|pw}\pend{}\pstart{}Vienna\oindex{Wien@\textbf{Wien}|pw}\pend{}\pstart{}{\pb}Hrn Gustav Schwarzkopf\pend{}\pstart{}Wien\oindex{Wien@\textbf{Wien}|pw}\pend{}\pstart{}I. Tiefer Graben 17\oindex{Wien@\textbf{Wien}!I., Innere Stadt@\textbf{I., Innere Stadt}!Tiefer Graben 17@\textbf{Tiefer Graben 17}|pw}\pend{}{\bigskip}
\pstart
           \noindent{}{\pb}\textcolor{gray}{\textbf{VENEZIA – Chiesa di S. Marco\oindex{Piazza San Marco@\textbf{Piazza San Marco}|pw} –}}\pend
           
\pstart
           \textcolor{gray}{\textbf{Torre dell’ Orologio\oindex{Torre dell’Orologio@\textbf{Torre dell’Orologio}|pw}.}}\pend
           \vspace{1em}
\pstart
           \raggedleft{}{\pb}Venezia\oindex{Venedig@\textbf{Venedig}|pw}{ }22. 4. 927\pend
           \vspace{0.5em}
\pstart
           Viele herzliche Grüße\pend
           
\pstart
           Ihr{\\[\baselineskip]}\spacefill\mbox{Arthur}\pend
           \leftskip=0em{}\selectlanguage{ngerman}\endnumbering\briefempfaengerindex{Schwarzkopf, Gustav@\textsc{Schwarzkopf, Gustav}!zzzSchnitzler, Arthur@\emph{von Arthur Schnitzler}!1927-04-221@{22. 4. 1927}|)be}\mylabel{L04468h}
\begin{anhang}
\end{anhang}\newcommand{\dateiname}{L04468}\newcommand{\titel}{Arthur Schnitzler an Gustav Schwarzkopf, 22. 4. 1927}\newcommand{\editorInnen}{Herausgegeben von Jahnke, SelmaMüller, Martin Anton}\input{../tex-inputs/latex-leseansicht-abspann}
